		
		
	
	\section{Demographic Survey} \label{Appendix:DemographicSurvey}
	\vspace{12pt}
	\noindent
	\textbf{1. Full Name: \hrulefill}
	
	\vspace{12pt}
	\noindent
	\textbf{2. Age: \hrulefill}
	
	\vspace{12pt}
	\noindent				
	\textbf{3. Education Level:}
	
	\begin{itemize}
		\item No formal education ( )
		\item Primary ( )
		\item Secondary ( )
		\item Higher education ( )
	\end{itemize}
	
	\noindent
	\textbf{4. Gender:}
	
	\begin{itemize}
		\item Male ( )
		\item Female ( )
		\item Other ( )
	\end{itemize}
	
	\noindent
	\textbf{5. How would you describe your experience monitoring your child's distractions while doing schoolwork?}
	
	\begin{itemize}
		\item Easy ( )
		\item Average ( )
		\item Challenging ( )
		\item Very challenging ( )
	\end{itemize}
	
	\noindent
	\textbf{6. How frequently do you monitor your child’s academic activities?}
	
	\begin{itemize}
		\item Never ( )
		\item Sometimes ( )
		\item Almost always ( )
		\item Always ( )
	\end{itemize}
	
	\noindent
	\textbf{7. What strategies have you implemented to improve your child’s concentration during schoolwork?}
	
	\vspace{12pt}
	\noindent
	\rule{\textwidth}{0.5pt}	
	\\
	\\
	\noindent
	\rule{\textwidth}{0.5pt}
	
	\clearpage
	
	\section{Questionnaires to Evaluate the Usability and Effectiveness of Torddis} \label{Appendix:Questionnaires}
		This appendix compiles the instruments used to assess both the usability and effectiveness of the Torddis system during its pilot implementation. The questionnaires were designed following validated models such as the System Usability Scale (SUS) and adapted to the educational and domestic context in which the prototype was deployed. The collected responses contribute significantly to understanding user satisfaction, perceived utility, and the potential pedagogical impact of the system. This section includes three types of questions: a Likert-scale questionnaire to measure usability, a set of open-ended questions to capture qualitative insights, and a final Likert-based instrument focused on evaluating the system’s effectiveness in real-world use.
	
		\subsection{Likert Scale Questionnaire} \label{Appendix:SUSLikertQuestionary}
			The Likert-scale questionnaire is based on the internationally recognized System Usability Scale (SUS), which provides a standardized framework to evaluate the usability of interactive systems. This instrument was selected due to its reliability and widespread application in human-computer interaction studies. The SUS allows capturing both positive and negative perceptions regarding the ease of use, complexity, integration, and learnability of the system. In the context of Torddis, the questionnaire was adapted to reflect specific aspects of its use by guardians in a home-based educational environment. The following items were presented to participants after a defined period of system interaction. Table \ref{tab:SUSLikertQuestions} lists the questions, and respondents were asked to indicate their agreement level using the following scale:
			
			\begin{enumerate}
				\item Strongly disagree
				\item Disagree
				\item Neutral (neither agree nor disagree)
				\item Agree
				\item Strongly agree
			\end{enumerate}
			
			\begin{table}[h]
				\caption{SUS questionnaire \label{tab:SUSLikertQuestions}}
				\centering
				\begin{tabularx}{\textwidth}{c X c c c c c }
					\toprule
					\textbf{No.} & \textbf{Questions} & 1 & 2 & 3 & 4 & 5 \\
					\midrule
					1 & I would like to use the Torddis system frequently. & & & & & \\
					2 & I found the system unnecessarily complex. & & & & & \\
					3 & I thought the system was easy to use. & & & & & \\
					4 & I would need technical or instructor support to use the system. & & & & & \\
					5 & I found the various system functions to be well integrated (forming a coherent whole). & & & & & \\
					6 & I thought the system had too many inconsistencies. & & & & & \\
					7 & I imagine most people would learn to use the system very quickly. & & & & & \\
					8 & I found the system very difficult to use. & & & & & \\
					9 & I felt very confident/comfortable using the system. & & & & & \\
					10 & I needed to learn many things before I could use the system. & & & & & \\
					\bottomrule
				\end{tabularx}
				\parbox{\textwidth}{\footnotesize 
					1: Strongly disagree. 2: Disagree. 3: Neutral (neither disagree nor agree). 4: Agree. 5: Strongly agree.
				}
			\end{table}
		
		\subsection{Open-Ended Questions from the SUS Questionnaire} \label{Appendix:SUSOpenQuestionarie}
			In addition to the structured Likert-scale items, a set of open-ended questions was included to capture more nuanced and personalized perspectives from the users. These questions allowed participants to express their opinions, preferences, and suggestions freely, thereby enriching the quantitative data with qualitative insights. The open-ended format is particularly valuable in educational technology research, as it reveals motivations, emotional responses, and contextual factors that may not be evident through closed questions. Table \ref{tab:OpenQuestion} presents the open-ended questions aimed at deepening the understanding of user experience with both the Torddis device and mobile application.
			
			\begin{table}[h]
				\centering
				\caption{Open-ended questions from the SUS questionnaire \label{tab:OpenQuestion}}
				\begin{tabularx}{\textwidth}{c X}
					\toprule
					\textbf{No.} & \textbf{Question}\\
					\midrule
					1 & In general, what is your opinion about the system?\\
					2 & Did you like the sounds and/or lights included in the Torddis device? Please answer yes or no and provide a reason.\\
					3 & Did you like the screen design of the Torddis mobile application? Please answer yes or no and provide a reason.\\
					4 & Is the way data is displayed about your child’s distraction monitoring in the Torddis mobile app appropriate? Why?\\
					5 & Do you believe this system could help improve your child’s concentration and keep you informed when they are distracted? Please answer yes or no and provide a reason.\\
					6 & Do you think there is anything that should be improved in the Torddis system (Device and Mobile App)? If so, what?\\
					7 & Would you be willing to continue using the Torddis system? Please answer yes or no and provide a reason.\\
					8 & Would you recommend this system to others interested in monitoring children’s distraction while doing schoolwork? Please answer yes or no and provide a reason.\\
					\bottomrule
				\end{tabularx}
			\end{table}
			
		\subsection{Questions to Assess the Effectiveness of the Prototype} \label{Appendix:LikertEffectivity}
			The following section presents a dedicated Likert-scale questionnaire focused on assessing the perceived effectiveness of the Torddis prototype in real educational contexts. These items were specifically designed to evaluate the functionality, relevance, and pedagogical contribution of the system after its application in home environments. Unlike the general usability assessment, this questionnaire targets outcomes such as attention improvement, self-regulation, system attractiveness, and user engagement. Responses to these questions provide empirical evidence regarding the system’s practical value, contributing to its validation as a meaningful educational technology. Table \ref{tab:EffectivityQuestions} presents the list of effectiveness-related items.
			
			\begin{table}[h]
				\centering
				\caption{Likert scale questions to assess system effectiveness \label{tab:EffectivityQuestions}}	
				\begin{tabularx}{\textwidth}{c X c c c c c }
					\toprule
					\textbf{No.} & \textbf{Questions} & 1 & 2 & 3 & 4 & 5 \\
					\midrule
					1 & The device was perceived as attractive by the child.  & & & & &  \\
					2 & The mobile application correctly recorded the events that occurred.  & & & & &  \\
					3 & The child did not perceive the presence of the device as intrusive.  & & & & &  \\
					4 & The device was used throughout the entire time dedicated to schoolwork.  & & & & &  \\
					5 & After discontinuing the use of the device, the child demonstrated improvement in self-regulation and concentration while completing tasks.  & & & & &  \\
					6 & Even though the device was used for a short period, it is expected to significantly contribute to task completion and academic achievement.  & & & & &  \\
					7 & You would be willing to continue using the device.  & & & & &  \\
					\bottomrule
				\end{tabularx}
				\parbox{\textwidth}{\footnotesize 
					1: Strongly disagree. 2: Disagree. 3: Neutral (neither disagree nor agree). 4: Agree. 5: Strongly agree.
				}
			\end{table}


